\begin{abstract}
A prerequisite for fixing issues in large repositories is code localization - the identification of relevant files, classes, and functions from the repository that require modification. 
While repository-level code localization was previously performed using embedding-based retrieval approaches like vector search, recent work has focused on developing LLM agents for this task, either as a standalone precursor to or interleaved with the code repair process.
Most previously published research studies on agentic code search have focused on providing the agent with complex tools such as repository graphs using static analysis tools, or vector indices.
In this paper, we demonstrate that, with a strong reinforcement learning recipe, a coding agent equipped with \emph{nothing more} than the standard Unix terminal, can be trained to achieve surprisingly strong results, beating all previously reported results for its model size, and sometimes approaching performace provided by closed models such as Claude Sonnet even when using specialized scaffolds.
In particular, we focus on the techniques used to re-purpose existing coding agent environments into environments for code search, reward design, and RL optimization techniques.\sonicomment{TODO: Summarize the best results here}
We release the resulting model family, \methodname, and all training code and data for the community to build upon.\footnote{\methodname is available at \url{https://github.com/OpenHands/codescout}}
% Effective localization requires navigating complex software repositories and rapidly acquiring contextual understanding. Despite this promise, most existing approaches are either closed-source or design complex scaffolds that are difficult to use and do not generalize across programming languages. In this work, we present an open-source recipe for training a dedicated code localization agent. We re-purpose established issue-resolution benchmarks (SWE-Smith and SWE-Bench) for localization tasks and integrate them with the OpenHands agent harness. Using this framework, we train \texttt{\methodname} (Code-Search OSS), a specialized code localization agent that achieves superior or competitive performance across SWE-Bench Lite, Verified, and Pro.\sonicomment{TODO: add concrete absolute gains here and a summarized variant of results}
\end{abstract}